\emph{Conventions.} Rotors, the bivectors $B_{k}$, the Haar measure and its moments $\int u^{\alpha}dU = \prod_{k}(\alpha_{k} - 1)!!/\prod_{j=0}^{|\alpha|/2 - 1}(4 + 2j)$ (zero unless every $\alpha_{k}$ is even) are those of Theorem \claimref{synthesis.weak-factor-rotor-lattice-strong-coupling-area-law}(i), (ii); the check multiplies rotors by a four-component product compared first with \texttt{GA.mul} on all sixteen basis pairs. Lattice weights, Haar measure and reflection positivity are those of Fact \claimref{synthesis.imported-wilson-lattice-gauge-theory-and-strong-coupling-confinement}.

\emph{Step 1 (the continuum copy; proved).} Write $q = P/D$ with $P = (1 - r^{2}) + 2\sum_{k}x_{k}B_{k}$ and $D = 1 + r^{2}$. Then $P\widetilde P = (1 - r^{2})^{2} + 4r^{2} = D^{2}$, so $q\widetilde q = 1$. The gauge transform of $A = 0$ by $q$ is $A_{i} = \widetilde q\,\partial_{i}q = \widetilde P(D\,\partial_{i}P - P\,\partial_{i}D)/D^{3}$, a bivector-valued rational field; at the origin $A_{1} = 2B_{1} \neq 0$. For the line element $\Omega^{2}$ times the flat one in three dimensions, the divergence of a covector field $A$ is $\Omega^{-3}\sum_{i}\partial_{i}(\Omega A_{i})$ (in index form, $\sqrt{g}\,g^{ij} = \Omega\,\delta_{ij}$, a parenthetical reading), so the Coulomb condition of the round three-sphere is $\sum_{i}\partial_{i}(\Omega A_{i}) = 0$. With $N_{i} = \widetilde P(D\,\partial_{i}P - P\,\partial_{i}D)$, $\Omega A_{i} = 2N_{i}/D^{4}$ and $D^{5}\sum_{i}\partial_{i}(N_{i}/D^{4}) = \sum_{i}(D\,\partial_{i}N_{i} - 4N_{i}\,\partial_{i}D)$, a polynomial valued in the even subalgebra of $\Cl(3,0)$; the check expands it with exact coefficients and finds it zero, while the flat analogue $\sum_{i}(D\,\partial_{i}N_{i} - 3N_{i}\,\partial_{i}D)$ is not zero. (Standard reading: $q$ is the identity map of the three-sphere onto the unit rotors, and $\widetilde q\,dq$ is its left-invariant Maurer-Cartan form, whose components are dual to Killing fields of the round metric and hence co-closed.) For an abelian field the copies of $A = 0$ are $A = d\theta$; the condition gives $\Delta_{S^{3}}\theta = 0$, and on a compact connected manifold $\int\theta\,\Delta\theta = -\int|d\theta|^{2}$ forces $d\theta = 0$. This is (i).

\emph{Step 2 (the lattice copy; computed).} Stationary points of $g \mapsto F_{U}(g)$ at $g = 1$ are the configurations whose lattice divergence $\sum_{\mu}(\langle U_{x,\mu}\rangle_{2} - \langle U_{x - \mu,\mu}\rangle_{2})$ of bivector parts vanishes at every site (first variation of $\langle(1 + \epsilon w_{x})U(1 - \epsilon w_{y})\rangle_{0}$), the lattice Landau gauge. Since $B_{3}^{4} = 1$, $g_{(x,y)} = B_{3}^{x}$ is periodic on four sites; the transformed links are $B_{3}^{x}B_{3}^{-(x+1)} = -B_{3}$ along the first axis and $1$ along the second. Both configurations have constant links and hence zero divergence. The second variation along $g_{x} = 1 + \epsilon w_{x} + \tfrac12\epsilon^{2}w_{x}^{2}$ is $\epsilon^{2}\sum_{\ell}\langle\tfrac12w_{x}^{2}U_{\ell} + \tfrac12U_{\ell}w_{y}^{2} - w_{x}U_{\ell}w_{y}\rangle_{0}$; the check builds its $24 \times 24$ matrix on the three bivector components at the eight sites by polarisation and computes the signature of the Faddeev-Popov form, minus twice that coefficient, by exact congruence: $(21, 0, 3)$ at $U = 1$, where it is the lattice Laplacian with the three global modes as kernel, and $(14, 2, 8)$ at the copy. A configuration whose Faddeev-Popov form has a negative mode is a saddle of $F$ on its orbit and lies outside the Gribov region. The gauge function uses only $B_{3}$, so the copy has a compact $U(1)$ analogue (links $-B_{3}$ in the group $a + bB_{3}$), where the Faddeev-Popov form at the copy is semi-definite; the negative modes come from the off-diagonal terms $\langle B_{a}B_{3}B_{b}\rangle_{0}$, $a, b \in \{1, 2\}$, which couple neighbouring sites and exist only for the non-abelian group. This is (ii).

\emph{Step 3 (the signed count; computed).} With the first rotor fixed and $g$ at the second site, $F(g) = \langle U_{1}\widetilde g\rangle_{0} + \langle gU_{2}\rangle_{0} = \langle g\widetilde{U_{1}}\rangle_{0} + \langle gU_{2}\rangle_{0} = \langle gM\rangle_{0}$, using $\langle X\rangle_{0} = \langle\widetilde X\rangle_{0}$. Along $g\exp(w)$ the first variation is $\langle gwM\rangle_{0} = \langle wMg\rangle_{0}$, zero for all bivectors $w$ exactly when $Mg$ is a scalar, that is $g = \lambda\widetilde M$ with $\lambda$ real; unit length gives $\lambda = \pm(M\widetilde M)^{-1/2} = \pm\tfrac58$. So there are exactly two stationary points. The second variation is $\langle g\,\tfrac12(w_{a}w_{b} + w_{b}w_{a})\,M\rangle_{0} = -\delta_{ab}\langle gM\rangle_{0} = \mp\tfrac85\,\delta_{ab}$, so the Faddeev-Popov forms are $\pm\tfrac85$ times the identity, of determinant signs $+1$ and $-1$; the check computes both matrices and their signatures. Their sum is $0 = \chi(S^{3})$, as the Poincar\'e-Hopf theorem requires for a Morse function on the three-sphere. Neuberger's theorem \cite{Neuberger1987}, as read through \cite{KalloniatisVonSmekalWilliams2005} (Fact \claimref{synthesis.imported-singer-gauge-slice-neuberger-lattice-brst-and-fradkin-shenker-complementarity}(ii)), is the same count for the gauge-fixing functional on any lattice and any compact group: the integral of a BRST-invariant observable over the gauge orbit with the Faddeev-Popov determinant (with sign) is a topological invariant proportional to the Euler characteristic of the gauge group, which vanishes for $SU(2)$ and for $U(1)$; so the gauge-fixed partition function is zero and gauge-fixed expectation values are $0/0$. This general statement is cited, not proved. This is (iii).

\emph{Step 4 (perimeter law at $\beta = 0$; proved).} The weight is $\prod_{\ell}\exp(\kappa\langle\widetilde{\varphi_{x}}U_{\ell}\varphi_{y}\rangle_{0})$ with Haar measure on links and on the unit rotors $\varphi$. The change of variables $U_{\ell} \mapsto \varphi_{x}U_{\ell}\widetilde{\varphi_{y}}$ preserves the Haar measure and $\langle U_{C}\rangle_{0}$ (Theorem \claimref{synthesis.weak-factor-rotor-lattice-strong-coupling-area-law}(iv), a gauge transformation by $g = \varphi$), and turns the weight into $\prod_{\ell}\exp(\kappa\langle U_{\ell}\rangle_{0})$ (unitary gauge; the check tests this by integrating over the site rotors with the Haar moments: the hopping term has the moments of $u_{0}$ up to the fourth, and the two-link loop $0 \to 1 \to 0$, integrated over both sites and both links without gauge fixing, has the expansion $\tfrac{\kappa^{2}}{16} - \tfrac{\kappa^{4}}{192} + O(\kappa^{6})$, which is $c(\kappa)^{2}$ to that order); the $\varphi$ integrals then factor out. The links are now independent, each with the conjugation-invariant weight $e^{\kappa u_{0}}$. Its mean $\bar U$ commutes with every rotor, so $\bar U = c(\kappa)$ is a scalar, and $\overline{\widetilde U} = c(\kappa)$ as well. For $C$ using $P$ distinct links once each, independence gives $\langle U_{C}\rangle = c(\kappa)^{P}$ and $\langle\langle U_{C}\rangle_{0}\rangle = c(\kappa)^{P}$, exactly. Expanding, $c = (\kappa m_{2} + \tfrac{\kappa^{3}}{6}m_{4} + \dots)/(1 + \tfrac{\kappa^{2}}{2}m_{2} + \dots)$ with $m_{2} = \int u_{0}^{2}dU = \tfrac14$ and $m_{4} = \tfrac18$, so $c = \tfrac{\kappa}{4} + (\tfrac{1}{48} - \tfrac{1}{32})\kappa^{3} + O(\kappa^{5}) = \tfrac{\kappa}{4} - \tfrac{\kappa^{3}}{96} + O(\kappa^{5})$. The check computes the moments, the gluing value $\int U u_{0}\,dU = \tfrac14$ and $\int U\,dU = 0$ exactly, using the 24 Hurwitz units after checking that they reproduce every Haar moment of degree at most three, and the leading coefficients $4^{-4}$ and $4^{-6}$ of the $1 \times 1$ and $1 \times 2$ loops. Since the Haar measure is symmetric under $u_{0} \mapsto -u_{0}$, the numerator of $c$ is $\int_{u_{0} > 0}u_{0}(e^{\kappa u_{0}} - e^{-\kappa u_{0}})\,dU > 0$, so $c(\kappa) > 0$ for $\kappa > 0$, the loop decays as $e^{-P\ln(1/c)}$, a perimeter law. This is (iv).

\emph{Step 5 (reading; argued).} (a) At small $\beta > 0$ and small $\kappa > 0$ the leading terms of a large $R \times T$ loop are $(\kappa/4)^{2(R + T)}$ and $(\beta/4)^{RT}$ (Theorem \claimref{synthesis.weak-factor-rotor-lattice-strong-coupling-area-law}(v)); the first dominates as the loop grows, so the area law of the pure gauge field does not survive the doublet. The full statement, that no phase boundary separates the Higgs and confinement regions for a fundamental Higgs field, is \cite{FradkinShenker1979} (Fact \claimref{synthesis.imported-singer-gauge-slice-neuberger-lattice-brst-and-fradkin-shenker-complementarity}(iii)). The series' Higgs is recorded as two fundamental doublets (\claimref{synthesis.linear-rung-field-has-a-neutral-normal-pair-in-the-scalon-sum}); that pattern is input here, not derived, and the frozen-length model is the standard proxy. The stored area-law theorem already says its area law is not evidence of confinement; the perimeter law is what is added. (b) The Krein premise is used by the BRST quantisation at the free level and in loops; Steps 1 to 3 show that its gauge fixing is not global beyond perturbation theory, and Gribov's ambiguity (Fact \claimref{synthesis.imported-brst-kugo-ojima-strocchi-and-gribov}(iv)) and Singer's theorem \cite{Singer1978} (Fact \claimref{synthesis.imported-singer-gauge-slice-neuberger-lattice-brst-and-fradkin-shenker-complementarity}(i)) are the general statements. Neither refutes the premise (its refutation conditions concern the physical space and gauge dependence of amplitudes), but positivity under (K4) beyond perturbation theory cannot be had from a lattice BRST construction as it stands. (c) As a reading, not adopted by any stored result: the gauge-invariant rotor lattice is a non-perturbative definition available for the weak factor (no colour lattice is stored); it has a positive transfer matrix by reflection positivity (Fact \claimref{synthesis.imported-wilson-lattice-gauge-theory-and-strong-coupling-confinement}(ii)) without gauge fixing, ghosts or an indefinite form. (d) Continuum confinement of pure colour, with the asymptotic freedom of \claimref{synthesis.one-loop-running-coefficients-from-spacetime-algebra-spin-traces}, is the open mass-gap question; nothing here bears on it.

\emph{Status.} (i) and (iv) proved; (ii) and (iii) computed exactly on the stated lattices, with the general signed count cited; the reading argued. All standard; recomputed in the rotors of the weak factor. Transport-absent.
